\documentclass[11pt]{article}
\usepackage[margin=1in]{geometry}
\usepackage{amsmath,amssymb,amsthm,mathtools}
\usepackage{enumitem}
\usepackage{microtype}
\input{glyphtounicode}
\pdfgentounicode=1
\pdfglyphtounicode{mu}{03BC}
\pdfglyphtounicode{bardbl}{2016}
\pdfglyphtounicode{parenleftbig}{0028}
\pdfglyphtounicode{parenrightbig}{0029}
\pdfglyphtounicode{parenleftbigg}{0028}
\pdfglyphtounicode{parenrightbigg}{0029}
\pdfglyphtounicode{parenlefttp}{239B}
\pdfglyphtounicode{parenleftbt}{239D}
\pdfglyphtounicode{parenrighttp}{239E}
\pdfglyphtounicode{parenrightbt}{23A0}
\pdfglyphtounicode{circleplusdisplay}{2A01}
\pdfglyphtounicode{summationdisplay}{2211}
\pdfglyphtounicode{productdisplay}{220F}
\usepackage[colorlinks=true,linkcolor=blue,citecolor=blue,urlcolor=blue]{hyperref}
\usepackage{accsupp}
\NewCommandCopy{\OriginalHookrightarrow}{\hookrightarrow}
\NewCommandCopy{\OriginalLongmapsto}{\longmapsto}
\NewCommandCopy{\OriginalUpbracefill}{\upbracefill}
\renewcommand{\hookrightarrow}{\mathrel{
  \BeginAccSupp{method=hex,unicode,ActualText=21AA}
  \OriginalHookrightarrow\EndAccSupp{}}}
\renewcommand{\longmapsto}{\mathrel{
  \BeginAccSupp{method=hex,unicode,ActualText=27FC}
  \OriginalLongmapsto\EndAccSupp{}}}
\renewcommand{\upbracefill}{
  \BeginAccSupp{method=hex,unicode,ActualText=23DF}
  \OriginalUpbracefill\EndAccSupp{}}
\hypersetup{pdftitle={Quadratic stabilization of the canonical Foulkes--Howe map},pdfauthor={OpenAI}}
\newtheorem{theorem}{Theorem}
\newtheorem{corollary}[theorem]{Corollary}
\newtheorem{proposition}[theorem]{Proposition}
\newtheorem{lemma}[theorem]{Lemma}
\theoremstyle{remark}
\newtheorem{remark}[theorem]{Remark}
\DeclareMathOperator{\Sym}{Sym}
\DeclareMathOperator{\Av}{Av}
\newcommand{\C}{\mathbb C}
\newcommand{\Sa}{\mathfrak S_a}
\newcommand{\D}[2]{D_{#1,#2}}
\newcommand{\unverified}[1]{\par\smallskip\noindent\fbox{\parbox{0.95\linewidth}{\textbf{Unverified step.} #1}}\par\smallskip}
\title{Quadratic stabilization of the canonical Foulkes--Howe map}
\author{Rewritten for readability from the OpenAI manuscript of the same title}
\date{September 25, 2026}
\begin{document}
\maketitle
\begin{abstract}
For every finite-dimensional complex vector space $V$, the canonical
Foulkes--Howe map $\Sym^b(\Sym^a V)\to\Sym^a(\Sym^b V)$ is surjective
whenever $a\ge2$ and $b\ge a(a-1)$. The bound is independent of $\dim V$.
\end{abstract}

\section{The result}\label{sec:introduction}

Throughout, $V$ is a finite-dimensional complex vector space, and products
of vectors are taken in the symmetric algebra
$\Sym V=\bigoplus_{j\ge0}\Sym^jV$.

\paragraph{Averaging.}
For a complex vector space $W$ and an integer $a\ge1$, the group $\Sa$
permutes the factors of $W^{\otimes a}$, and
\begin{equation}\label{eq:average}
 \Av(w_1\otimes\cdots\otimes w_a)
 =\frac1{a!}\sum_{\pi\in\Sa}
 w_{\pi(1)}\otimes\cdots\otimes w_{\pi(a)}
\end{equation}
defines a linear projection of $W^{\otimes a}$ onto the invariants
$(W^{\otimes a})^{\Sa}$. The right side is symmetric in the $w_i$, so
$w_1\cdots w_a\mapsto\Av(w_1\otimes\cdots\otimes w_a)$ is a well-defined map
$\Sym^aW\to(W^{\otimes a})^{\Sa}$. It is an isomorphism: its inverse is the
quotient map $W^{\otimes a}\to\Sym^aW$ restricted to invariants.

\paragraph{The map.}
Fix $a\ge1$ and put
\[
 R_j=\bigl((\Sym^j V)^{\otimes a}\bigr)^{\Sa}\quad(j\ge0),\qquad
 R=\bigoplus_{j\ge0}R_j.
\]
The algebra $(\Sym V)^{\otimes a}$, with factorwise multiplication, is
commutative, and $\Sa$ acts on it by algebra automorphisms; hence $R$ is a
graded commutative subalgebra. Averaging identifies $R_1=\Sym^aV$ and
$R_b=\Sym^a(\Sym^bV)$. For $b\ge0$, multiplying $b$ elements of $R_1$ gives a
linear map $\Sym^bR_1\to R_b$. Under these identifications it is the
\emph{canonical Foulkes--Howe map}
\begin{equation}\label{eq:map}
 \mu_{a,b,V}:\Sym^b(\Sym^a V)\longrightarrow\Sym^a(\Sym^b V).
\end{equation}
Its image is $A_b$, the degree-$b$ part of the subalgebra $A\subseteq R$
generated by $R_1$. (In positive dimension, Raicu--Sam--Weyman
\cite[Section~2]{RSW} formulate this map as
$\Sym^b(D^aV)\to D^a(\Sym^bV)$, with divided powers realized as invariant
tensors; averaging transports it to \eqref{eq:map}.)

\begin{theorem}\label{thm:main}
For every integer $a\ge2$, every integer $b\ge a(a-1)$, and every
finite-dimensional complex vector space $V$, the map
$\mu_{a,b,V}$ in \eqref{eq:map} is surjective. Equivalently, $A_b=R_b$.
\end{theorem}

For $a=1$ the map is the identity in every degree. Corollary~\ref{cor:comparison} deduces
a $\mathrm{GL}(V)$-equivariant embedding
$\Sym^a(\Sym^bV)\hookrightarrow\Sym^b(\Sym^aV)$ in the same range.

\paragraph{Proof at a glance.}
The idea is to dualize surjectivity into a vanishing statement for symmetric
multilinear forms, and then to deduce that a polynomial is zero from the
vanishing of a chain of its directional derivatives.
\begin{enumerate}[nosep]
\item \emph{Criterion} (Section~\ref{sec:map}). $A_b$ is spanned by the
tensors $P^{\otimes a}$ with $P=u_1\cdots u_b$ a product of $b$ vectors
(Lemmas~\ref{lem:powers},~\ref{lem:image}). Dually
(Proposition~\ref{prop:criterion}), $\mu_{a,b,V}$ is surjective if and only if
every symmetric $a$-linear form $T$ on $\Sym^bV$ with $T(P,\ldots,P)=0$ for
all such $P$ is zero.
\item \emph{Injectivity} (Section~\ref{sec:differential}). The derivative
$\D{Y}{U}=\sum_kY_k\,\partial/\partial U_k$ is injective on polynomials whose
$U$-degree exceeds their $Y$-degree (Lemma~\ref{lem:shift}). So a polynomial
is zero if a chain of such derivatives of it is zero and the degree inequality
holds before each derivative (Corollary~\ref{cor:chain}).
\item \emph{Vanishing} (Proposition~\ref{prop:vanishing},
Section~\ref{sec:vanishing}). If $b\ge a(a-1)$, every $T$ as in item~1 is
zero. The proof is by induction on $a$; write $r=a-1$ and $m=b-r\ge r^2$.
 \begin{itemize}[nosep]
 \item Step 1: test $T$ on $P=(x^r+ct^r)Q$, which is a product of $b$ vectors
 because binary forms split over $\C$. The coefficient of $c$ gives
 $T(x^rQ,\ldots,x^rQ,t^rQ)=0$.
 \item Step 2: $m$ derivatives turn $t^b$ into a multiple of $t^rQ$, so
 Corollary~\ref{cor:chain} gives $T(x^rQ,\ldots,x^rQ,t^b)=0$.
 \item Step 3: the induction hypothesis, applied to the first $r$ slots,
 replaces the common $Q$ by arbitrary $Q_1,\ldots,Q_r\in\Sym^mV$.
 \item Step 4: $r^2$ derivatives turn $T(v_1^b,\ldots,v_r^b,t^b)$ into a
 multiple of $T(x^rv_1^m,\ldots,x^rv_r^m,t^b)=0$. They all raise the degree
 of the same block $x$, up to $r^2-1$, while the source degrees stay
 $\ge m+1$. This accumulation produces the quadratic bound.
 \end{itemize}
\end{enumerate}
Dependencies: Lemma~\ref{lem:powers} $\to$ Lemma~\ref{lem:image} $\to$
Proposition~\ref{prop:criterion}; Lemma~\ref{lem:shift} $\to$
Corollary~\ref{cor:chain} $\to$ Steps 2 and 4; Steps
$1\to2\to3\to4$; Propositions~\ref{prop:criterion} and~\ref{prop:vanishing}
$\to$ Theorem~\ref{thm:main} $\to$ Corollary~\ref{cor:comparison}.

\paragraph{Context.}
The affine Chow variety of decomposable $a$-forms consists of products of $a$
linear forms in $\Sym^a(V^*)$; its coordinate ring identifies with $A$, whose
normalization is $R$, with $b$ the coordinate degree
\cite[Section~2, Theorem~2.3]{RSW}; see also
\cite[Section~7.6, Propositions~7.11--7.12]{Landsberg}. Brion proved eventual
surjectivity \cite{Brion1993} and then an effective bound involving both $a$
and $\dim V$ \cite[Theorem~3.3]{Brion}. Raicu, Sam, and Weyman bound the
regularity of the Segre ring and modules of covariants, including the
normalization \cite[Theorem~3.9]{RSW}; this does not compute the top degree
of the quotient $R/A$, which controls stabilization.
Theorem~\ref{thm:main} gives the bound $a(a-1)$ in every dimension and
answers positively the request for a polynomial bound in
\cite[Problem~7.19]{Landsberg} (there $d=b$, $n=a$, $w=\dim V$; the bound is
$d\ge n(n-1)$, including $w=n$). The bound is not asserted to be sharp.
The question differs from the Foulkes conjecture \cite{Foulkes1950}, which
asks, for $a\le b$, for some $\mathrm{GL}(V)$-equivariant injection
$\Sym^a(\Sym^bV)\hookrightarrow\Sym^b(\Sym^aV)$
\cite[Conjecture~2.10]{RSW}. Surjectivity of the canonical map implies such
an injection (Corollary~\ref{cor:comparison}), but an injection need not come
from the canonical map: at $a=b$ the two representations coincide, yet the
canonical map has nonzero kernel at $(a,b)=(5,5)$ and $(6,6)$ in suitable
dimensions \cite[arXiv version, Theorem~6(b),(c) and Remark~7]{CIM}, both
outside our range. Successive weight shifts also occur in Ikenmeyer's proof
\cite[Section~5]{Ikenmeyer} of McKay's propagation theorem \cite{McKay}:
for the opposite canonical maps $\Sym^a(\Sym^bV)\to\Sym^b(\Sym^aV)$ with
$\dim V\ge b$, injectivity at $(a,B)$, $B\ge a$, propagates to all $(a,b)$
with $b\ge B$ \cite[Theorem~1.3 and Section~3]{Ikenmeyer}.

\section{The canonical image and diagonal vanishing}\label{sec:map}

Fix integers $a,b\ge1$ and put $E=\Sym^bV$. A \emph{product of $b$ vectors}
is an element $P=u_1\cdots u_b\in E$ with $u_i\in V$; the vectors may repeat
or be zero.

\begin{lemma}\label{lem:powers}
For every integer $k\ge1$, the pure powers $u^k$ $(u\in V)$ span $\Sym^kV$.
\end{lemma}

This is used in Lemma~\ref{lem:image} ($k=a$) and at the end of the proof of
Proposition~\ref{prop:vanishing} ($k=b$).

\begin{proof}
We claim the polarization identity
\[
 k!\,u_1\cdots u_k
 =\sum_{J\subseteq\{1,\ldots,k\}}(-1)^{k-|J|}
   \Bigl(\sum_{j\in J}u_j\Bigr)^k .
\]
Treat the $u_i$ as formal commuting variables and expand. A degree-$k$
monomial with support $S\subseteq\{1,\ldots,k\}$ occurs in
$(\sum_{j\in J}u_j)^k$ exactly when $J\supseteq S$, with a multinomial
coefficient $\gamma$ independent of $J$. Its total coefficient is
$\gamma\sum_{J\supseteq S}(-1)^{k-|J|}=\gamma(1-1)^{k-|S|}$. This is $0$
unless $S=\{1,\ldots,k\}$, in which case the monomial is $u_1\cdots u_k$ and
$\gamma=k!$. Substituting vectors gives the identity in $\Sym V$. Products
span $\Sym^kV$ and $k!\ne0$, so pure powers span.
\end{proof}

\begin{lemma}\label{lem:image}
$A_b=\operatorname{span}_{\C}\{P^{\otimes a}:P\text{ a product of }b\text{
vectors}\}$.
\end{lemma}

This is used only in Proposition~\ref{prop:criterion}.

\begin{proof}
(i) Averaging sends $u^a\in\Sym^aV$ to $z(u):=u^{\otimes a}$. Averaging maps
$\Sym^aV$ onto $R_1$, so by Lemma~\ref{lem:powers} with $k=a$ the elements
$z(u)$ span $R_1$.
(ii) $A_b$ is spanned by products of $b$ elements of $R_1$. Multiplication is
multilinear, so $A_b$ is spanned by the products $z(u_1)\cdots z(u_b)$.
(iii) Factorwise multiplication gives
$z(u_1)\cdots z(u_b)=(u_1\cdots u_b)^{\otimes a}$.
\end{proof}

\begin{proposition}[Diagonal criterion]\label{prop:criterion}
The map $\mu_{a,b,V}$ is surjective if and only if every symmetric $a$-linear
form $T$ on $E=\Sym^bV$ satisfying $T(P,\ldots,P)=0$ for all products $P$ of
$b$ vectors is zero.
\end{proposition}

The proof of Theorem~\ref{thm:main} uses the ``if'' direction.

\begin{proof}
For $\lambda\in R_b^*$ define the symmetric $a$-linear form
$T_\lambda(w_1,\ldots,w_a)=\lambda\bigl(\Av(w_1\otimes\cdots\otimes
w_a)\bigr)$ on $E$.

(i) \emph{$\lambda\mapsto T_\lambda$ is a bijection from $R_b^*$ onto the
symmetric $a$-linear forms on $E$.} It is injective because $\Av$ maps
$E^{\otimes a}$ onto $R_b$, so averages of simple tensors span $R_b$. It is
surjective: given a symmetric $T$, let $\widetilde T:E^{\otimes a}\to\C$ be
its linearization and $\lambda=\widetilde T|_{R_b}$. Then
$T_\lambda(w_1,\ldots,w_a)=\frac1{a!}\sum_\pi
T(w_{\pi(1)},\ldots,w_{\pi(a)})=T(w_1,\ldots,w_a)$ by symmetry.

(ii) For a product $P$, the tensor $P^{\otimes a}$ is invariant, so
$T_\lambda(P,\ldots,P)=\lambda(P^{\otimes a})$. By Lemma~\ref{lem:image},
$\lambda$ vanishes on $A_b$ if and only if $T_\lambda(P,\ldots,P)=0$ for all
products $P$.

(iii) $R_b$ is finite-dimensional, so $A_b=R_b$ if and only if the only
$\lambda\in R_b^*$ vanishing on $A_b$ is $0$. By (i) and (ii) this is the
stated condition on forms, and $\operatorname{im}\mu_{a,b,V}=A_b$.
\end{proof}

\section{Differential injectivity}\label{sec:differential}

\paragraph{Conventions.}
Assume $n=\dim V>0$ and fix a basis $e_1,\ldots,e_n$. A \emph{vector
variable} is a block $U=(U_1,\ldots,U_n)$ of independent scalar variables,
standing for $\sum_kU_ke_k$. An expression built from vector variables by
products in $\Sym V$ and a multilinear form expands, by multilinearity, to a
polynomial in the coordinates, homogeneous in each block. A polynomial over
$\C$ that vanishes at every point is zero (induction on the number of
variables, from the univariate root bound). Hence such a polynomial is zero
once the expression vanishes for all vector values.

For two blocks $U,Y$ define
\[
 \D{Y}{U}=\sum_{k=1}^n Y_k\frac{\partial}{\partial U_k}.
\]
It lowers the $U$-degree by one, raises the $Y$-degree by one, and leaves the
degree in every other block unchanged. By the chain rule,
$\D{Y}{U}h=\frac{d}{ds}h(U+sY)\big|_{s=0}$, with the other blocks held fixed.
Let $L$ be a linear form on $\Sym^{k+j}V$ and $g\in\Sym^jV$, neither
involving $U$ (they may depend polynomially on other blocks, including $Y$).
Since $(U+sY)^kg=U^kg+sk\,U^{k-1}Yg+O(s^2)$ in $\Sym V$,
\begin{equation}\label{eq:insert}
 \D{Y}{U}\,L(U^kg)=k\,L(U^{k-1}Y\,g)\qquad(k\ge1).
\end{equation}

\begin{lemma}\label{lem:shift}
Let $p,q$ be nonnegative integers. On polynomials homogeneous of degree $p$
in $U$ and degree $q$ in $Y$, the operator $\D{Y}{U}$ is injective if
$p>q$. The polynomials may depend on any finite collection of additional
variables independent of $U$ and $Y$.
\end{lemma}

This is the positive-weight injectivity in $\mathfrak{sl}_2$ weight strings;
compare \cite[Section~11.1, Claim~11.4 and equation~(11.5)]{FultonHarris}.
Ikenmeyer's argument \cite[Section~5, Claim~5.1]{Ikenmeyer} decomposes
tensors according to untouched positions before applying a two-coordinate
weight shift. The proof below is a direct inner-product argument with all
additional variables present. The condition $p>q$ cannot be dropped: for
$p=q=1$, $h=U_1Y_2-U_2Y_1$ satisfies $\D{Y}{U}h=0$. (Renaming: the source
calls the operator $\D{U}{Y}$ below $E$; we write $D'$, since $E$ denotes
$\Sym^bV$.)

\begin{proof}
Let $X$ list all coordinates, including the additional variables. Use the
Fischer (apolar) Hermitian inner product
\cite[Section~2, equation~(8) and Proposition~3]{AldazRender}: distinct
monomials are orthogonal and
$\|X^\alpha\|^2=\alpha!:=\prod_k\alpha_k!$. It is positive definite.

\emph{Adjoints.} For each coordinate $X_j$, $\partial_{X_j}$ is adjoint to
multiplication by $X_j$: on monomials,
$\langle\partial_{X_j}X^\alpha,X^\beta\rangle$ and
$\langle X^\alpha,X_jX^\beta\rangle$ both equal $\alpha!$ if
$\alpha=\beta+e_j$ and $0$ otherwise (using
$\alpha_j\cdot(\alpha-e_j)!=\alpha!$). Put $D=\D{Y}{U}$ and $D'=\D{U}{Y}$.
Then $(Y_k\partial_{U_k})^*=(\partial_{U_k})^*(Y_k)^*=U_k\partial_{Y_k}$, so
$D^*=D'$ and $D'^*=D$.

\emph{Commutator.} By the Leibniz rule,
\[
 U_j\partial_{Y_j}\,Y_k\partial_{U_k}
 =\delta_{jk}U_j\partial_{U_k}+U_jY_k\partial_{Y_j}\partial_{U_k},
 \qquad
 Y_k\partial_{U_k}\,U_j\partial_{Y_j}
 =\delta_{jk}Y_k\partial_{Y_j}+Y_kU_j\partial_{U_k}\partial_{Y_j}.
\]
Subtracting and summing over $j,k$,
\[
 D'D-DD'=N_U-N_Y,
 \qquad
 N_U=\sum_k U_k\partial_{U_k},\quad
 N_Y=\sum_k Y_k\partial_{Y_k}.
\]
On a monomial, $N_U$ multiplies by the $U$-degree and $N_Y$ by the
$Y$-degree.

\emph{Conclusion.} For $h$ of bidegree $(p,q)$ in $(U,Y)$,
\begin{equation}\label{eq:norm}
 \|Dh\|^2-\|D'h\|^2
 =\langle D'Dh,h\rangle-\langle DD'h,h\rangle
 =\langle(N_U-N_Y)h,h\rangle=(p-q)\|h\|^2.
\end{equation}
If $Dh=0$ and $p>q$, then $(p-q)\|h\|^2=-\|D'h\|^2\le0$, so $h=0$.
\end{proof}

\begin{corollary}\label{cor:chain}
Let $G_0$ be a polynomial homogeneous in each block, and let
$G_i=\D{Y^{(i)}}{U^{(i)}}G_{i-1}$ for $1\le i\le N$. Suppose that $G_N=0$ and
that, for every $i$, the degree of $G_{i-1}$ in $U^{(i)}$ exceeds its degree
in $Y^{(i)}$. Then $G_0=0$.
\end{corollary}

\begin{proof}
Each $G_{i-1}$ is homogeneous in each block. Lemma~\ref{lem:shift} applied to
$G_{N-1}$ gives $G_{N-1}=0$; repeat downward to $G_0$.
\end{proof}

Corollary~\ref{cor:chain} is used in Steps 2 and 4 below.

\section{Vanishing on products}\label{sec:vanishing}

\begin{proposition}\label{prop:vanishing}
Let $V$ be a finite-dimensional complex vector space, and let $a\ge1$ and
$b\ge1$ be integers satisfying $b\ge a(a-1)$.
If a symmetric $a$-linear form $T$ on $\Sym^b V$ satisfies
\begin{equation}\label{eq:diagonal}
 T(P,\ldots,P)=0
 \qquad\text{for every product $P$ of $b$ vectors},
\end{equation}
then $T=0$.
\end{proposition}

\emph{Plan.} We induct on $a$. In the inductive step, Step~1 produces a
vanishing with a common factor $Q$ in every slot, Step~2 removes $Q$ from the
last slot, Step~3 applies the induction hypothesis to the other $r=a-1$
slots, and Step~4 removes the common power $x^r$ from those slots.

\begin{proof}
If $V=0$ then $\Sym^bV=0$ and $T=0$. So let $\dim V>0$ and use the
conventions of Section~\ref{sec:differential}. We induct on $a$; the
statement for a given $a$ covers every admissible $b$ and every $V$.

\emph{Base case $a=1$.} $T$ is a linear form vanishing on all products of
$b$ vectors, and these span $\Sym^bV$.

\emph{Inductive step.} Let $a\ge2$, assume the statement for $a-1$, and put
\begin{equation}\label{eq:rm}
 r=a-1\ge1,\qquad m=b-r.
\end{equation}
Since $b\ge a(a-1)=r(r+1)$,
\begin{equation}\label{eq:range}
 m\ge r^2\ge1.
\end{equation}

\paragraph{Step 1: a split binary form.}
\emph{Claim: for all $x,t\in V$ and every product $Q=y_1\cdots y_m$ of $m$
vectors,}
\begin{equation}\label{eq:binary}
 T(\underbrace{x^rQ,\ldots,x^rQ}_{r\text{ slots}},t^rQ)=0.
\end{equation}
\begin{enumerate}[label=(\alph*),nosep]
\item \emph{For every $c\in\C$, $P_c=(x^r+ct^r)Q$ is a product of $b$
vectors.} Choose $\alpha\in\C$ with $\alpha^r=-c$ and let
$\zeta=e^{2\pi i/r}$. From $\prod_{j=0}^{r-1}(X-\beta\zeta^j)=X^r-\beta^r$
with $\beta=\alpha Y$,
\[
 X^r+cY^r=\prod_{j=0}^{r-1}(X-\alpha\zeta^jY)\quad\text{in }\C[X,Y].
\]
The algebra homomorphism $\C[X,Y]\to\Sym V$, $X\mapsto x$, $Y\mapsto t$
(which needs no independence of $x,t$) writes $x^r+ct^r$ as a product of $r$
vectors, and $r+m=b$.
\item By (a) and \eqref{eq:diagonal}, $T(P_c,\ldots,P_c)=0$ for all $c$.
Since $P_c=x^rQ+c\,t^rQ$, multilinearity and symmetry give
\[
 T(P_c,\ldots,P_c)=\sum_{k=0}^{a}\binom ak c^k\,
 T(\underbrace{x^rQ,\ldots,x^rQ}_{a-k},\underbrace{t^rQ,\ldots,t^rQ}_{k}).
\]
A polynomial in $c$ vanishing on $\C$ has zero coefficients. The coefficient
of $c$ is $a\,T(x^rQ,\ldots,x^rQ,t^rQ)$ with $a-1=r$ slots $x^rQ$, and
$a\ne0$.
\end{enumerate}

\paragraph{Step 2: freeing the last slot.}
\emph{Claim: for all $x,t\in V$ and every product $Q$ of $m$ vectors,
$T(x^rQ,\ldots,x^rQ,t^b)=0$ ($r$ slots $x^rQ$).}

Fix $x\in V$. Let $t,y_1,\ldots,y_m$ be vector variables, $Q=y_1\cdots y_m$,
and
\[
 G_0=H_x=T(\underbrace{x^rQ,\ldots,x^rQ}_{r\text{ slots}},t^b),
 \qquad G_i=\D{y_i}{t}G_{i-1}\quad(1\le i\le m).
\]
\begin{enumerate}[label=(\alph*),nosep]
\item \emph{Formula.} For $0\le i\le m$,
\begin{equation}\label{eq:first-composition}
 G_i=\frac{b!}{(b-i)!}\,
 T(x^rQ,\ldots,x^rQ,\;t^{b-i}y_1\cdots y_i).
\end{equation}
This follows by induction on $i$ from \eqref{eq:insert} with $U=t$, $Y=y_i$,
$k=b-i+1$, $g=y_1\cdots y_{i-1}$ and $L=T(x^rQ,\ldots,x^rQ,\,\cdot\,)$; the
first $r$ slots do not involve $t$.
\item \emph{$G_m=0$.} Since $b-m=r$, $G_m$ is a multiple of
$T(x^rQ,\ldots,x^rQ,t^rQ)$. This vanishes for all vector values by
\eqref{eq:binary}, hence as a polynomial.
\item \emph{Degrees.} By \eqref{eq:first-composition}, $G_{i-1}$ has
$t$-degree $b-i+1$ and $y_i$-degree $r$: $y_i$ occurs once in each of the $r$
copies of $Q$ and not yet in the last slot. The difference is
$(b-i+1)-r=m-i+1\ge1$.
\item Corollary~\ref{cor:chain} gives $H_x=G_0=0$, which is the claim.
\end{enumerate}

\paragraph{Step 3: induction on the remaining slots.}
\emph{Claim: for all $x,t\in V$,}
\begin{equation}\label{eq:inductive-vanishing}
 T(x^rQ_1,\ldots,x^rQ_r,t^b)=0
 \qquad(Q_1,\ldots,Q_r\in\Sym^m V).
\end{equation}
Fix $x,t$. Multiplication by $x^r$ is linear $\Sym^mV\to\Sym^bV$, so
$S_{x,t}(Q_1,\ldots,Q_r):=T(x^rQ_1,\ldots,x^rQ_r,t^b)$ is a symmetric
$r$-linear form on $\Sym^mV$. By Step~2, $S_{x,t}(Q,\ldots,Q)=0$ for every
product $Q$ of $m$ vectors. The induction hypothesis applies with $(r,m)$ in
place of $(a,b)$, because $r\ge1$, $m\ge1$ and $m\ge r^2\ge r(r-1)$ by
\eqref{eq:range}. Hence $S_{x,t}=0$.

\paragraph{Step 4: removing the common power.}
\emph{Claim: $T(v_1^b,\ldots,v_r^b,t^b)=0$ for all
$v_1,\ldots,v_r,t\in V$.}

Fix $t\in V$. Let $v_1,\ldots,v_r$ and $x$ be vector variables (so $x$ is now
a variable block) and
\[
 f=T(v_1^b,\ldots,v_r^b,t^b),
\]
of degree $b$ in each $v_i$ and degree $0$ in $x$. Apply $r^2$ operators in
this order: $\D{x}{v_1}$ $r$ times, then $\D{x}{v_2}$ $r$ times, and so on up
to $\D{x}{v_r}$.
\begin{enumerate}[label=(\alph*),nosep]
\item \emph{Result.} Apply \eqref{eq:insert} $r$ times in slot $i$, with
$U=v_i$, $Y=x$ and $g$ a power of $x$; the other slots do not involve $v_i$.
This replaces $v_i^b$ by
$b(b-1)\cdots(b-r+1)\,x^rv_i^{b-r}=\frac{b!}{m!}x^rv_i^m$. Hence
\begin{equation}\label{eq:second-composition}
 \D{x}{v_r}^{\,r}\cdots\D{x}{v_1}^{\,r}f
 =\left(\frac{b!}{m!}\right)^r
 T(x^rv_1^m,\ldots,x^rv_r^m,t^b)=0.
\end{equation}
The last equality holds for all vector values by
\eqref{eq:inductive-vanishing} with $Q_i=v_i^m$, hence as a polynomial
identity.
\item \emph{Degrees.} Before the $j$-th of the $r^2$ operators, the
$x$-degree is $j-1\le r^2-1$. The source block $v_i$ of that operator has
been differentiated at most $r-1$ times, so its degree is at least
$b-(r-1)=m+1\ge r^2+1$ by \eqref{eq:range}. So the source degree exceeds the
$x$-degree.
\item Corollary~\ref{cor:chain} gives $f=0$.
\end{enumerate}

\paragraph{Conclusion.}
By Step~4, $T$ vanishes on every $a$-tuple of pure $b$-th powers. These span
$\Sym^bV$ (Lemma~\ref{lem:powers}, $k=b$), so $T=0$ by multilinearity. This
completes the induction.
\end{proof}

\section{Proof of the theorem; representation comparison}\label{sec:comparison}

\begin{proof}[Proof of Theorem~\ref{thm:main}]
Let $a\ge2$ and $b\ge a(a-1)$; then $b\ge1$.
Proposition~\ref{prop:vanishing} shows that every symmetric $a$-linear form
on $\Sym^bV$ whose diagonal vanishes on products of $b$ vectors is zero.
Proposition~\ref{prop:criterion} then gives surjectivity of $\mu_{a,b,V}$.
\end{proof}

\begin{corollary}\label{cor:comparison}
For every finite-dimensional complex vector space $V$, every integer
$a\ge2$, and every integer $b\ge a(a-1)$, there exists a
$\mathrm{GL}(V)$-equivariant embedding
\[
 \Sym^a(\Sym^b V)\hookrightarrow\Sym^b(\Sym^a V).
\]
In particular, this gives the sixth symmetric-power comparison for $b\ge30$.
\end{corollary}

\begin{proof}
Averaging and multiplication commute with the action of $\mathrm{GL}(V)$, so
the surjection $\mu_{a,b,V}$ of Theorem~\ref{thm:main} is equivariant.
Finite-dimensional polynomial representations over $\C$ are completely
reducible \cite[Chapter~I, Appendix~A, Sections~7--8]{Macdonald}. Choose an
invariant complement $M$ to $\ker\mu_{a,b,V}$. Then $\mu_{a,b,V}|_M$ is an
equivariant isomorphism onto $\Sym^a(\Sym^bV)$, and its inverse followed by
the inclusion $M\subseteq\Sym^b(\Sym^aV)$ is the embedding. This proves
existence only; the embedding is not identified with a prescribed canonical
map.
\end{proof}

\begin{remark}
Lemma~\ref{lem:shift} holds for any number of coordinates and additional
variables, and the degree comparisons involve only $a$ and $b$. With the
finite-dimensional annihilator argument, this gives the same bound in every
finite dimension, without a stable-dimension reduction. The proof uses
positivity of the Fischer inner product, splitting of binary forms over
$\C$, and the nonzero characteristic-zero scalars ($a$, $k!$, $b!/(b-i)!$,
$b!/m!$) in the coefficient and differentiation identities.
\end{remark}

\section{Verification checklist}

\begin{enumerate}[nosep]
\item \textbf{Duality} (Proposition~\ref{prop:criterion}, (i)--(ii)): the
restricted linearization $\lambda$ satisfies $T_\lambda=T$, and
$T_\lambda(P,\ldots,P)=\lambda(P^{\otimes a})$ with $P^{\otimes a}$ spanning
$A_b$ (Lemma~\ref{lem:image}).
\item \textbf{Injectivity lemma} (Lemma~\ref{lem:shift}): the adjointness
$D^*=D'$, the commutator $D'D-DD'=N_U-N_Y$, and the sign in
\eqref{eq:norm}.
\item \textbf{Step 1(a)}: $(x^r+ct^r)Q$ is a product of exactly $r+m=b$
vectors for every $c$ and all $x,t$, including dependent ones. \textbf{Step
1(b)}: the coefficient of $c$ is $a\,T(x^rQ,\ldots,x^rQ,t^rQ)$.
\item \textbf{Step 2(c)}: before the $i$-th derivative, the $y_i$-degree is
exactly $r$ and the $t$-degree is $b-i+1$, so the difference is
$m-i+1\ge1$.
\item \textbf{Step 3}: the induction hypothesis is applied to $(r,m)$ and
needs $m\ge1$ and $m\ge r(r-1)$; both follow from \eqref{eq:range}.
\item \textbf{Step 4(b)}: all $r^2$ derivatives raise the degree of the same
block $x$ (at most $r^2-1$ beforehand), while each source degree is at least
$m+1\ge r^2+1$. Check also the constant $(b!/m!)^r\ne0$ in
\eqref{eq:second-composition}.
\item \textbf{Polynomial versus pointwise vanishing}: Steps 2(b) and 4(a)
pass from vanishing at all vector values to a polynomial identity
(Conventions, Section~\ref{sec:differential}), which
Corollary~\ref{cor:chain} requires.
\end{enumerate}

\begin{thebibliography}{99}
\bibitem{AldazRender}
J.~M.~Aldaz and H.~Render,
\href{https://doi.org/10.1007/s13324-023-00844-4}{\emph{A Fischer type decomposition theorem from the apolar inner product}},
Analysis and Mathematical Physics \textbf{13} (2023), article~91.

\bibitem{Brion1993}
M.~Brion,
\href{https://doi.org/10.1007/BF03026558}{\emph{Stable properties of plethysm: on two conjectures of Foulkes}},
Manuscripta Mathematica \textbf{80} (1993), no.~4, 347--371.

\bibitem{Brion}
M.~Brion,
\href{https://smf.emath.fr/publications/sur-certains-modules-gradues-associes-aux-produits-symetriques}{\emph{Sur certains modules gradu\'es associ\'es aux produits sym\'etriques}},
in \emph{Alg\`ebre non commutative, groupes quantiques et invariants}
(Reims, 1995), S\'eminaires et Congr\`es \textbf{2},
Soci\'et\'e Math\'ematique de France, Paris, 1997, 157--183.

\bibitem{CIM}
M.-W.~Cheung, C.~Ikenmeyer, and S.~Mkrtchyan,
\href{https://doi.org/10.1016/j.jsc.2016.09.002}{\emph{Symmetrizing tableaux and the 5th case of the Foulkes conjecture}},
Journal of Symbolic Computation \textbf{80} (2017), part~3, 833--843.
\href{https://arxiv.org/abs/1509.03944v1}{arXiv:1509.03944v1}.

\bibitem{Foulkes1950}
H.~O.~Foulkes,
\href{https://doi.org/10.1112/jlms/s1-25.3.205}{\emph{Concomitants of the quintic and sextic up to degree four in the coefficients of the ground form}},
Journal of the London Mathematical Society \textbf{25} (1950), no.~3,
205--209.

\bibitem{FultonHarris}
W.~Fulton and J.~Harris,
\href{https://doi.org/10.1007/978-1-4612-0979-9}{\emph{Representation Theory: A First Course}},
Graduate Texts in Mathematics \textbf{129}, Springer-Verlag, New York, 1991.

\bibitem{Ikenmeyer}
C.~Ikenmeyer,
\emph{On McKay's propagation theorem for the Foulkes conjecture},
\href{https://arxiv.org/abs/1509.04957v1}{arXiv:1509.04957v1}, 2015.

\bibitem{Landsberg}
J.~M.~Landsberg,
\href{https://doi.org/10.1007/s11565-014-0202-7}{\emph{Geometric complexity theory: an introduction for geometers}},
Annali dell'Universit\`a di Ferrara \textbf{61} (2015), no.~1, 65--117.
\href{https://arxiv.org/abs/1305.7387}{arXiv:1305.7387}.
Locators refer to the published version.

\bibitem{Macdonald}
I.~G.~Macdonald,
\href{https://doi.org/10.1093/oso/9780198534891.001.0001}{\emph{Symmetric Functions and Hall Polynomials}},
second edition, Oxford Mathematical Monographs, Clarendon Press, Oxford, 1995.

\bibitem{McKay}
T.~McKay,
\href{https://doi.org/10.1016/j.jalgebra.2007.12.003}{\emph{On plethysm conjectures of Stanley and Foulkes}},
Journal of Algebra \textbf{319} (2008), no.~5, 2050--2071.

\bibitem{RSW}
C.~Raicu, S.~V.~Sam, and J.~Weyman,
\href{https://doi.org/10.1007/s10013-021-00527-2}{\emph{On some modules supported in the Chow variety}},
Vietnam Journal of Mathematics \textbf{50} (2022), no.~2, 501--521.
\href{https://arxiv.org/abs/2108.10910v1}{arXiv:2108.10910v1}.
Locators refer to this preprint version.
\end{thebibliography}

\end{document}
